\paragraph{Coordinates on a short arc.}\label{par:short-wach-frame}
\begin{proposition}[Common coordinates for short-arc comparisons]\label{prop:short-arc-coordinates}
For a short arc, use the same integral vertical and
horizontal frames and the corresponding cyclotomic levels.
For an arc ending at \(c\) of length \(a=a_c\), the offset \(0\le j<a\)
has size \(u_j=s_{c-j}\) and cyclotomic factor \(P_j=\phi^j(X+\ell)\). The horizontal and vertical comparisons use the same positive pole and negative zero factors
\[
 D=\prod_{j\ {\rm positive}}P_j^{u_j},\qquad V=\prod_{j\ {\rm negative}}P_j^{u_j+S}.
\]
The two ordinary lifting rules use these coordinates. In the single-letter edge cases the
positive rule applies to every input, and the negative rule
uses the empty pole set and the initial constant lift.
\end{proposition}

\begin{proof}
\proofstep{Fix the same cut coordinate and ordered cyclotomic levels for both rules.}
Write \(a=a_c,\ c'=c-a,\ G=w+S,\ u_j=s_{c-j}\), and
\(P_j=\phi^j(q)\), so \(P_0=q=X+\ell\).
Let \(\mathcal R=O[[X]]=O[[q]]\). All ideals of fractions below
are in the component at \(c\).
Call an offset positive or negative according to its large Hodge
jumps. This sign is used even at the first letter of small size.
The full denominator for transport over \(a\) steps and the two
smaller products determined by the jumps are
\[
 \Lambda=P_{a-1}^{G}\prod_{j<a-1}P_j^h,\qquad
 D=\prod_{j\ {\rm positive}}P_j^{u_j},\qquad
 V=\prod_{j\ {\rm negative}}P_j^{u_j+S}.
\]
Choose a basis of the tame-invariant Wach lattice at \(c'\).
The unique formal \(\gamma\)-fixed basis with the same constants is
its horizontal basis.

Write \(U(X)=I+\sum_{r>0}U_rX^r\) for the change
to the horizontal basis, and let \(F_r(U_1,\ldots,U_{r-1})\)
be the degree-\(r\) residual computed with \(U_r=0\).
The recursion is
\[
       (\omega_X^r-1)U_r=-F_r(U_1,\ldots,U_{r-1}).
\]
Its new denominator has order \(1+v_\ell(r)\) at step \(r\).
Thus through that step the denominators divide
\(\ell^{r+v_\ell(r!)}\). This also holds for the inverse change
of basis, by the factorial inequality for products of coefficients.
Use the basis truncated before degree \(M\). These are coordinates
on the given lattice map: a map of Wach lattices acts on horizontal
columns by its matrix of constants at \(c'\), since both sides solve
the same equation with the same constant. Frobenius transports these
constants to the crystalline space at \(c\).

Here are details on the precision of these coordinates. Put
\[
 P(T)=\langle\Phi^aN_{c'}(T)\rangle_{\mathcal R}.
\]

Represent \(v\in\Lambda^{-1}P(T)\) modulo \(X\Lambda P(T)\)
in the transported horizontal basis, allowing powers of \(\ell\)
to clear coefficients in this comparison. Equivalently multiply
by \(\Lambda\) and compare modulo \(X\Lambda^2 P(T)\).
The tails beyond the truncation have a factor
\(\phi^a(X)^M\), divisible by \(X\Lambda^2\) since
\((\ell-1)M>2h\). The finite coefficients, and the inverse matrix,
can be transported integrally at cost \(\sigma\). Indeed on
the first step one reduces \(\ell^\sigma\phi(X)^r/\ell^{r+v_\ell(r!)}\)
modulo \(X^Mq^{2G}\).

The term \(X^r q^{(\ell-1)r}\) pays the
denominator after this precision, by the binomial division of
Paragraph~\ref{par:wach-cut-correction} with one forward step and
precision \(2G\). More explicitly, for \(k=r+v_\ell(r!)-\sigma>0\)
one needs
\[
 (\ell-1)r\ge 2G+k+M-r-1.
\]
Then modulo \(\ell^k\)
the binomial expansion of \((X+\ell)^{(\ell-1)r-2G}\) is divisible
by \(X^{M-r}\). If \(k>0\), the coefficient bound gives
\(r>(\ell-1)\sigma/\ell\). Together with the definition of
\(\sigma\) this yields the preceding power inequality.

In
the binomial expansion the terms of \(\ell\)-degree \(\ge k\)
vanish modulo \(\ell^k\); every other term has \(X\)-degree
at least \((\ell-1)r-2G-k+1\ge M-r\). Thus the termwise
division is legitimate at the first precision. The next
\(a-1\) Frobenius maps give
\[
             \phi^{a-1}\bigl((X^Mq^{2G})\bigr)
                                     \subseteq (X\Lambda^2)
\]
by the same level bounds defining \(\Lambda\). Thus the transported matrices are represented by
integral lifts of \(\ell^\sigma\phi^a(U)\) and
\(\ell^\sigma\phi^a(U^{-1})\) modulo \(X\Lambda^2\)
in the Frobenius basis. If \(A,B\) denote these two integral lifts in the displayed order, their common precision gives
\[
 AB\equiv BA\equiv\ell^{2\sigma}I\pmod{X\Lambda^2}.
\]


For clarity, the coefficient test in these coordinates is as
follows. If a constant direction \(d\) at the end of the arc
belongs to \(\mathcal F_j^{-\alpha_j}\) for each transported
filtration, then \(d\) with coefficient in
\((\prod P_j^{\alpha_j})\mathcal R\) passes the Wach test at
each level. Sums of such terms after scalar extension give the
rational test; the bounded integral numerators are required as
well.

To see the rational assertion, separate the levels by the
Chinese remainder theorem for their distinguished polynomials
and take a basis of the filtration at each level. For a basis vector \(d\) of jump \(u\), the corresponding generator and its Laurent order are
\[
 t^{-u}d,\qquad \operatorname{ord}_t(t^{-u})=-u,
\]
as in
Paragraph~\ref{par:wach-level-models}. Multiplying the other
level factors to sufficiently high orders realizes each of these
finite jets without imposing further conditions there.
Divisibility of bounded numerators by a product of these
polynomials descends to \(\mathcal R\) by monic division.

Let \(u\) denote a coordinate vector satisfying all the
local tests and let \(n\) be a vector in \(N_c\).
The forward comparison and its inverse give, in their
respective quotients by the transported precision,
\[
   f_{\rm hor}(u)=[\ell^\sigma u]_{\rm Wach},
   \qquad g_{\rm hor}(n)=[\ell^\sigma n]_{\rm coordinates},
   \qquad g_{\rm hor}f_{\rm hor}=\ell^{2\sigma}.
\]
Thus an admissible coordinate vector produces a Wach class,
and the restricted inverse sends a Wach vector back into the
tested coordinate module. Indeed the transported forward
and inverse matrices above give integral lifts with the bounded
denominator, and changes by \(X\Lambda P(T)\) lie in \(XN_c\)
by the height bounds.

For an admissible vector one tests
\((\gamma-1)/X\) in the same manner. This gives an integral
vector with \((\gamma-1)v\in XN_c\), since its horizontal
columns are fixed to the indicated precision. A product of the
\(P_j\) clears its denominators in \(N_c\). These statements
apply to subarcs beginning at \(c'\) with their own product
\(\Lambda\). \label{cor:inert-short-frame-1}

We will lift a coordinate \(f\) on a target horizontal column
\(d\) which belongs to \(\mathcal F_j^{u_j}\) at every positive
level. Its scalar hypotheses are
\begin{equation}\label{eq:short-scalar}
 Df\in\mathcal R,\qquad
 D(\ell^{k_j}f-x_j)\in P_j^{u_j+S}\mathcal R
 \quad(j\ {\rm negative})
\end{equation}
for some \(k_j\ge0,\ x_j\in O\). A lift at cost \(e\) represents
\(\ell^e f d\) by the image of an admissible source vector up to
\(XN_c(T_2)\) and the coordinate precision. There are two ways
to construct it.
\begin{enumerate}
\item If \(Df\in V\mathcal R\), then \(fd\) passes the target
Wach test. The poles at the positive levels are allowed by the
filtration of \(d\); the zeros at the negative levels accommodate
the lowest possible jump. The coordinate comparison and scalar image bound at \(c\) use successive multipliers
\[
 \ell^\sigma\ell^R=\ell^{R+\sigma}.
\]
Tame averaging preserves the lifted image, so no extra exponent enters this case.
\item Suppose the poles are confined to a subset of the positive
levels and an integral source direction maps to \(\ell^t d\) and
lies in \(\mathcal F_j^{u_j}\) there. At each omitted positive
level also require that \(f\) is constant modulo \(P_j^\eta\)
after clearing scalars and the retained denominator. Then the
coefficient \((\gamma f-f)/X\) times this source direction passes
the level tests. For each level factor, write \(\gamma P_j/P_j=1+Xh_j\) with \(h_j\in\mathcal R\). Then
\[
 (\gamma P_j-P_j)/X=h_jP_j.
\]

For every \(n\in\mathbb Z\) the ratio \((1+Xh_j)^n\) belongs to \(1+X\mathcal R\). Thus
\[
 (\gamma P_j^n-P_j^n)/X\in P_j^n\mathcal R ,
\]
including for negative powers. The constant disappears
at each omitted positive or negative level. Equation~\eqref{eq:short-scalar}
therefore supplies the required zeros at negative levels, and
\(\eta\) supplies them at omitted positive levels. The retained
positive poles pass by the source filtration. The admissible-vector
construction gives the lift at \(t+\sigma\). For an integral constant, transport the constant term of the initial target basis lift at \(c'\). Its coefficient differences vanish, so its cost is at most \(R+\sigma\).
\end{enumerate}

In the second rule the congruences use the denominator for the
retained positive levels. Costs of two summands may be padded
to their maximum, by linearity of the source map and of the
coordinate comparison. For a single positive letter the first
rule applies to every \(f\). For a single negative letter use
the second with empty pole set and the initial constant lift.
\end{proof}

